\begin{figure}[htbp]
\centering
\begin{tikzpicture}[
    node distance = 0.9cm and 1.2cm,
    font = \footnotesize,
    startstop/.style = {rectangle, rounded corners=10pt, minimum width=2.4cm, minimum height=0.7cm, text centered, draw=black!80, fill=gray!15},
    process/.style   = {rectangle, minimum width=3.4cm, minimum height=0.8cm, text centered, draw=blue!70!black, fill=blue!5, rounded corners=2pt, align=center},
    decision/.style  = {diamond, aspect=2, minimum width=2.8cm, text centered, draw=orange!80!black, fill=orange!10, align=center, inner sep=1pt},
    io/.style        = {trapezium, trapezium left angle=70, trapezium right angle=110, minimum width=3.0cm, minimum height=0.7cm, text centered, draw=teal!70!black, fill=teal!5, align=center},
    arrow/.style     = {-Stealth, thick, draw=gray!80}
]

    % Các khối chức năng
    \node (start) [startstop] {Bắt đầu};
    \node (fetch) [process, below=0.6cm of start] {Gửi HTTP Request đến API Camera\\{\scriptsize (\texttt{Get\_frame.py}, \texttt{cv2.imdecode})}};
    \node (time)  [process, below=0.6cm of fetch] {Ghi nhận Timestamp máy nhận ảnh};
    \node (check) [decision, below=0.6cm of time] {\texttt{frame.tobytes()} \\ trùng ảnh trước?};
    \node (yolo)  [process, below=0.7cm of check] {Chạy YOLO nhận diện phương tiện\\\& Tính điểm WTI};
    \node (buffer)[decision, below=0.6cm of yolo] {Đủ lô 5 mẫu\\hoặc ngoại lệ?};
    \node (write) [io, below=0.6cm of buffer] {Ghi dữ liệu vào CSV\\{\scriptsize (Batch flush / \texttt{finally})}};
    \node (loop)  [startstop, below=0.6cm of write] {Tiếp tục chu kỳ tiếp theo};

    % Đường nối và rẽ nhánh
    \draw [arrow] (start) -- (fetch);
    \draw [arrow] (fetch) -- (time);
    \draw [arrow] (time) -- (check);
    
    % Rẽ nhánh: Nếu trùng ảnh thì bỏ qua, quay lại lấy ảnh mới
    \draw [arrow] (check) -- node[right] {Sai (ảnh mới)} (yolo);
    \draw [arrow] (check.east) -- ++(1.4,0) node[above, pos=0.4] {Đúng (trùng)} |- (fetch.east);

    % Ghi dữ liệu
    \draw [arrow] (yolo) -- (buffer);
    \draw [arrow] (buffer) -- node[right] {Đúng} (write);
    \draw [arrow] (write) -- (loop);

    % Rẽ nhánh: Chưa đủ 5 mẫu thì gom tiếp
    \draw [arrow] (buffer.west) -- ++(-1.1,0) node[above, pos=0.4] {Chưa} |- (loop.west);

\end{tikzpicture}
\caption{Lưu đồ giải thuật quy trình thu thập ảnh, kiểm tra trùng lặp và ghi dữ liệu của mô-đun \texttt{get\_data.py}.}
\label{fig:get_data_flowchart}
\end{figure}